\source{Xu et al., ACL 2026, pp.122--146. DOI: 10.18653/v1/2026.acl-short.12}
\papersection{StructMem}{面向长程对话的结构化记忆}{hycx233}
\source{原文 §1--2；附录 A.8，表8，p.132。案例证据简化为中文}
\begin{frame}{长程对话中的证据组织}
\textbf{问题：}Caroline 和 Melanie 什么时候一起参加过 Pride festival？\par
\vspace{3mm}
\small
\begin{tabularx}{\linewidth}{@{}p{29mm}Y@{}}
\toprule
2023年8月的回忆 & 两人在交流中提到去年的活动经历，以及与朋友共同参与的记忆。\\[1mm]
回答需要的证据 & 将人物互动与活动经历关联起来，并由来源时间把“去年”换算为2022年。\\\bottomrule
\end{tabularx}\par
\vspace{4mm}
\textbf{三种组织方式}\par\vspace{1mm}
\begin{tabularx}{\linewidth}{@{}p{23mm}YY@{}}
\toprule
方法 & 保存形式 & 需要处理的问题\\\midrule
平面记忆 & 独立条目＋向量 & 检索后补齐分散的事件语境\\
图记忆 & 实体、关系与图结构 & 实体消歧、关系去重与维护\\
StructMem & 时间绑定的条目＋整合摘要 & 在写入时组织相关事件\\\bottomrule
\end{tabularx}\par
\aside{关键词相似能找到相关片段，完整回答还需要人物、事件与时间之间的联系。}
\note{建议45秒。\par
先看论文中的例子。问题问两个人什么时候一起参加过活动。对话中既有各自的经历，也有对对方的回应，还有“去年”这样的相对时间。模型需要把这些信息放在一起，才能得到共同参与关系和正确年份。
\par
常见的平面记忆把内容存成独立条目，检索很方便，但事件语境容易分散。图记忆显式保存关系，同时需要实体消歧、关系去重等操作。StructMem选择保留带时间的自然语言条目，再生成跨事件摘要，把部分证据组织工作提前到记忆写入阶段。}
\end{frame}

\source{原文图2、§3，pp.123--124；流程据原文重绘}
\begin{frame}{StructMem 的两层记忆}
\begin{center}
\begin{tikzpicture}[>=Latex,every node/.style={font=\small},flowstep/.style={draw=accent,text width=35mm,minimum height=10mm,align=center,inner sep=3pt}]
\node[flowstep] (a) {事实与关系抽取};
\node[flowstep,right=7mm of a] (b) {时间锚定};
\node[flowstep,right=7mm of b] (c) {事件缓存};
\node[flowstep,below=9mm of c] (d) {历史检索与事件重建};
\node[flowstep,left=7mm of d] (e) {跨事件整合};
\node[flowstep,left=7mm of e] (f) {条目＋摘要问答};
\draw[->] (a)--(b);\draw[->] (b)--(c);\draw[->] (c)--(d);\draw[->] (d)--(e);\draw[->] (e)--(f);
\end{tikzpicture}
\end{center}
\vspace{2mm}
\small
\begin{tabularx}{\linewidth}{@{}p{25mm}YY@{}}
\toprule
记忆层次 & 保存内容 & 回答时的作用\\\midrule
事件级记忆 & 事实、互动关系、来源时间 & 提供可核对的具体细节\\
跨事件记忆 & 当前事件与相关历史的整合摘要 & 提供跨越多次交流的联系\\\bottomrule
\end{tabularx}\par
\aside{两层共同保留。核心流程由提示调用与向量检索组成，无需训练新的语言模型。}
\note{建议45秒。\par
方法分为两层。第一层处理一段对话内部的信息：模型分别抽取事实和互动关系，把它们绑定到来源时间，再放入事件缓存。第二层处理不同事件之间的联系：系统检索相关历史，恢复事件上下文，然后生成整合摘要。
\par
回答问题时，两类记忆一起使用。事件条目保存具体细节，摘要提供已经整理过的联系。这也解释了结构化体现在哪里：条目通过时间绑定，历史通过摘要连接，形成了有层次的记忆。整个流程依靠提示调用和向量检索，不需要训练新的语言模型。}
\end{frame}

\source{原文 §3.1，式(1)--(2)；图4--7、13，pp.133--136、139}
\begin{frame}{双视角抽取与时间锚定}
\small
\textbf{对话示例：}Alice 分享研究生录取和机器学习方向，Bob 祝贺她、询问计划并表达支持。\par
\vspace{3mm}
\begin{tabularx}{\linewidth}{@{}p{25mm}Y@{}}
\toprule
事实视角 $\Phi_i$ & Alice 被研究生项目录取，计划研究机器学习。\\[1mm]
关系视角 $\Psi_i$ & Bob 为 Alice 感到高兴，支持她的选择并关注研究计划。\\\bottomrule
\end{tabularx}\par
\vspace{2mm}
\[
  \text{记忆条目}=\langle x, e_x, \tau_i\rangle,
  \qquad x\in\Phi_i\cup\Psi_i
\]
$x$：自然语言内容\quad $e_x$：检索向量\quad $\tau_i$：来源话语时间\par
\vspace{3mm}
\textbf{时间示例：}2023年8月说“去年参加过活动”。\par\vspace{1mm}
来源时间保留为2023年8月，回答时将事件年份换算为\key{2022年}。\par
\aside{抽取结果附带 source\_id；事件重建按相同来源时间聚合事实与关系条目。}
\note{建议50秒。\par
双视角分别回答两个问题：发生了什么，人们怎样回应和影响彼此。以论文的提示示例来说，Alice被录取、研究机器学习属于事实；Bob的祝贺、支持和兴趣属于互动关系。只保留前一类信息，就可能丢掉后续推理需要的人物联系。
\par
每个条目还保存检索向量和来源时间。这里要区分说话的时间与事件发生的时间。2023年说“去年参加活动”，原话时间仍然是2023年，活动年份则在回答时换算成2022年。抽取结果中的source id用于追踪来源，正文的事件重建规则则按相同时间锚补齐相关条目。}
\end{frame}

\source{原文 §3.2，式(3)--(5)；A.3；图8，pp.124、129、137}
\begin{frame}{历史检索与跨事件整合}
\small
\begin{enumerate}
\setlength{\itemsep}{2.5mm}
\item \textbf{积累当前事件：}条目按对话时间排序，以1小时窗口触发整合。
\item \textbf{检索历史种子：}拼接缓存文本并编码，按余弦相似度取 top-15。
\item \textbf{恢复事件语境：}取回每个种子同时间锚的其他事实与关系条目。
\item \textbf{生成整合摘要：}当前缓存与恢复后的历史事件共同进入整合模型。
\end{enumerate}
\vspace{1mm}
\[
E_{\tau}(x^*)=\{x'\in M\mid \tau(x')=\tau(x^*)\}
\]
\textbf{摘要要求：}200--350个英文词，保留人物、具体时间和行动，组织有依据的事件联系。\par
\vspace{3mm}
\textbf{关键区别：}15是历史种子数，展开后的上下文可能包含更多条目。\par
\aside{1小时指对话时间窗口。整合摘要与事件条目共同保留。}
\note{建议60秒。\par
跨事件整合可以拆成四步。首先，系统积累当前条目，并按对话时间排序，实验把窗口阈值设为一小时。然后把缓存文本拼接起来，编码成一个查询，用余弦相似度检索十五条历史种子。
\par
第三步是方法里的关键。系统不会只把这十五条内容交给模型，还会取回与种子共享时间锚的其他条目。这样，命中一个事实时，也能补回同一语境中的互动关系。最后，当前缓存与恢复后的历史事件一起生成摘要。
\par
提示要求摘要控制在两百到三百五十个英文词，保留具体时间、人物和行动，并说明有依据的联系。因此，检索负责找到入口，事件重建补足语境，整合负责组织跨事件信息。}
\end{frame}

\source{原文 A.3；图13，pp.129、139；案例参考表8，p.132}
\begin{frame}{回答阶段的证据组合}
\small
\begin{tabularx}{\linewidth}{@{}p{31mm}Yp{45mm}@{}}
\toprule
检索内容 & 提供的信息 & 回答中的用途\\\midrule
\key{60条事件记忆} & 具体事实、人物互动、来源时间 & 核对人物身份与细节\\[1mm]
\key{5条整合摘要} & 已整理的相关历史与事件联系 & 联合多段证据\\\bottomrule
\end{tabularx}\par
\vspace{4mm}
\textbf{回答提示中的处理规则}\par\vspace{2mm}
\begin{itemize}
\setlength{\itemsep}{1.5mm}
\item 同时检查两位说话人的记忆，区分说话人与记忆中提到的人物。
\item 依据来源日期换算“去年”“两个月前”等时间表达。
\item 信息冲突时优先较新的记忆，最终输出简短、具体的答案。
\end{itemize}
\aside{构建阶段的15条种子与回答阶段的60＋5分别控制不同环节。}
\note{建议40秒。\par
在回答阶段，系统同时检索六十条事件记忆和五条整合摘要。条目帮助核对具体人物、事实和时间，摘要提供跨事件联系。回答提示要求检查两位说话人的记忆，避免把提到的人物与说话人混淆，并把相对时间换算成明确日期。
\par
回到开头的问题，模型可以把活动事实、共同参与的交流语境和2023年的参照时间结合起来，得到2022年这个答案。这里六十加五控制回答上下文，前一页的十五条种子控制记忆构建，两组参数的作用不同。}
\end{frame}

\source{原文 §4.1；A.2--A.4、表3，pp.125、128--129}
\begin{frame}{实验设置与比较范围}
\small
\begin{columns}[T,onlytextwidth]
\column{.56\textwidth}
\textbf{LoCoMo 长期对话问答}\par\vspace{2mm}
10段对话，平均588轮、16,618 Token。\par\vspace{2mm}
\begin{tabularx}{\linewidth}{@{}lrY@{}}
\toprule 任务 & 数量 & 需要的能力\\\midrule
单跳 & 841 & 定位单一事实\\
多跳 & 282 & 连接多条证据\\
时间 & 321 & 推理时间关系\\
开放域 & 96 & 对话结合背景知识\\\midrule
合计 & 1,540 & 单跳占54.6\%\\\bottomrule
\end{tabularx}
\column{.40\textwidth}
\textbf{模型与指标}\par\vspace{2mm}
生成：gpt-4o-mini\par
向量：text-embedding-3-small\par\vspace{3mm}
正确率：LLM评判回答\par
构建成本：Token、调用数、耗时\par\vspace{3mm}
主要对比：全文上下文、RAG、图记忆及长期记忆系统。
\end{columns}
\vspace{3mm}
\textbf{读表方式：}总分结合分类结果看；检索预算各异，主表反映整套系统表现。
\note{建议45秒。\par
实验使用LoCoMo的十段长期对话，平均每段五百八十八轮，约一万六千六百个Token。论文评测四类任务，共一千五百四十题。单跳考察事实定位，多跳需要连接证据，时间任务考察时间关系，开放域还需要背景知识。
\par
基础模型使用gpt-4o-mini，向量模型也保持统一。回答正确与否由模型评判，成本统计记忆构建阶段。读表时有两个重点：单跳题占一半以上，总分需要结合分类成绩理解；各基线的检索预算不同，所以主表展示的是整体系统效果。}
\end{frame}

\source{原文表1，p.125。粗体为所列方法中最高值；差值单位为百分点}
\begin{frame}{总分提升与能力差异}
\vspace{2mm}
\small\renewcommand{\arraystretch}{1.25}
\begin{tabularx}{\linewidth}{@{}Yrrrrr@{}}
\toprule 方法 & 总分 & 单跳 & 多跳 & 时间 & 开放域\\\midrule
FullContext & 73.83 & \textbf{86.56} & 68.79 & 50.16 & 56.25\\
Zep & 75.14 & 79.79 & \textbf{74.11} & 67.71 & \textbf{66.04}\\
Memobase & 75.78 & 77.17 & 70.92 & \textbf{85.05} & 46.88\\
StructMem & \textbf{76.82} & 81.09 & 68.77 & 81.62 & 46.88\\\bottomrule
\end{tabularx}\par
\vspace{4mm}
\begin{columns}[T,onlytextwidth]
\column{.48\textwidth}
\textbf{主要收益}\par\vspace{2mm}
总分比Memobase高\key{1.04}。\par\vspace{1mm}
相对FullContext，时间推理提高\key{31.46}，总分提高2.99。
\column{.48\textwidth}
\textbf{能力取舍}\par\vspace{2mm}
相对FullContext，单跳下降5.47，多跳接近。\par\vspace{1mm}
Zep多跳、开放域更高；Memobase时间任务更高。
\end{columns}
\aside{时间相关任务受益明显；跨事件整合的具体作用，需要进一步看组件消融。}
\note{建议60秒。\par
主表中，StructMem的总准确率是百分之七十六点八二，比Memobase高一点零四个百分点。对比把全文直接放入上下文的方法，总分提高二点九九个百分点，其中时间推理从五十点一六提高到八十一点六二，变化最明显。
\par
但分类表现也说明了代价。单跳比全文上下文低五点四七个百分点，多跳几乎相同。全文保留了全部原始细节，直接事实查询仍然有优势。其他记忆系统也有擅长的任务：Zep的多跳和开放域更高，Memobase的时间分数更高。
\par
我的理解是，StructMem在细节保留和跨事件组织之间取得了较好的整体平衡。至于这种收益有多少来自整合模块，要结合后面的消融实验判断。}
\end{frame}

\source{原文表1、图3(b)，pp.125--126。M为百万Token，k为千Token}
\begin{frame}{构建成本与运行时间}
\small
\begin{columns}[T,onlytextwidth]
\column{.57\textwidth}
\textbf{完整实验的构建统计}\par\vspace{2mm}
\renewcommand{\arraystretch}{1.3}
\begin{tabularx}{\linewidth}{@{}Yrrr@{}}
\toprule 方法 & Token(M) & 调用数 & 时间(s)\\\midrule
MiniRAG & 10.103 & 2,508 & 2,566\\
StructMem & 1.937 & 1,056 & 22,854\\
Mem0g & 35.825 & 53,514 & 115,670\\\bottomrule
\end{tabularx}\par\vspace{3mm}
相对Mem0g：Token减少\key{94.6\%}，调用减少\key{98.0\%}，耗时减少80.2\%。
\column{.39\textwidth}
\textbf{同框架，首段对话}\par\vspace{2mm}
平面记忆\hfill77.6k\par
StructMem\hfill226.1k\par
图记忆\hfill402.3k\par\vspace{3mm}
StructMem约为平面记忆的2.91倍，比图记忆少43.8\%。
\end{columns}
\vspace{4mm}
\textbf{成本来自哪里：}双视角抽取和摘要增加调用；批量整合减少图维护的多阶段开销。\par
\aside{MiniRAG运行更快。调用量、Token消耗与墙钟时间需要分别评价。}
\note{建议50秒。\par
效率结果分两组看。左边是完整实验的构建统计。相对Mem0g，StructMem的Token消耗减少百分之九十四点六，调用数减少百分之九十八。不过MiniRAG实际运行更快，所以调用少、Token少和耗时短需要分别看。
\par
右边是同一框架、首段对话的比较。平面记忆最省，StructMem约为它的二点九一倍，又比图记忆少百分之四十三点八。这个结果说明，双视角抽取和摘要确实需要额外成本，而批量组织事件可以减少实体、关系维护的多阶段开销。论文体现的是组织能力与构建代价之间的取舍。}
\end{frame}

\source{原文表2，p.125。最后一行为加入跨事件整合后的百分点变化}
\begin{frame}{跨事件整合的消融结果}
\small\renewcommand{\arraystretch}{1.3}
\begin{columns}[T,onlytextwidth]
\column{.66\textwidth}
\begin{tabularx}{\linewidth}{@{}Yrrrr@{}}
\toprule 配置 & 单跳 & 多跳 & 时间 & 开放域\\\midrule
Flat Memory & 78.83 & 66.31 & 78.50 & 46.88\\
Graph Memory & 80.50 & 66.67 & 76.64 & 48.96\\
去掉跨事件整合 & 80.86 & 66.31 & 79.44 & 46.88\\
StructMem & 81.09 & 68.77 & 81.62 & 46.88\\\midrule
整合带来的变化 & +0.23 & \textbf{+2.46} & \textbf{+2.18} & 0.00\\\bottomrule
\end{tabularx}\par
\column{.30\textwidth}
\textbf{结果解释}\par\vspace{3mm}
多跳与时间任务需要联合分散证据，受益更大。\par\vspace{3mm}
单跳变化较小，开放域未改善。\par\vspace{3mm}
符合整合摘要以连接对话证据为主的功能。
\end{columns}
\aside{进一步实验可分别去掉关系抽取和时间锚定，定位每个设计的贡献。}
\note{建议50秒。\par
为了看清整合模块的作用，论文保持框架不变，移除跨事件整合进行比较。加回这个模块以后，多跳提高二点四六个百分点，时间推理提高二点一八；单跳只增加零点二三，开放域没有变化。
\par
这个分布与方法的目标吻合。多跳和时间任务都需要连接分散证据，整理好的跨事件摘要能提供帮助。单跳通常已经有直接事实可查，而开放域还依赖背景知识，所以受益较少。
\par
这组消融支持跨事件整合的价值。进一步研究时，可以分别去掉关系抽取和时间锚定，并控制检索Token预算，看看收益具体来自哪些环节。}
\end{frame}

\source{原文图8、表6--7，pp.130--131、137。错误关联率由模型分别评判}
\begin{frame}{记忆内容的忠实性}
\small
\textbf{事件条目：}三个评判模型报告的平均幻觉率为\key{2.36\%}。\par
\vspace{3mm}
\textbf{整合约束：}引用历史时标注具体时间，保留人物和行动细节，建立具体、有依据的联系。\par
\vspace{3mm}
\renewcommand{\arraystretch}{1.25}
\begin{tabularx}{\linewidth}{@{}Yrr@{}}
\toprule 评判模型 & 有约束的错误关联率 & 移除约束后\\\midrule
GPT & 0.61\% & 7.45\%\\
Qwen & 3.41\% & 20.00\%\\
DeepSeek & 3.63\% & 15.14\%\\\bottomrule
\end{tabularx}\par
\vspace{3mm}
\textbf{观察：}三个评判模型均报告错误关联增加，提示中保留时间与具体联系有实际作用。\par
\aside{条目幻觉与摘要关联采用不同统计单位；这些数值来自模型审核。}
\note{建议40秒。\par
生成式记忆还有一个问题：整理关系时会不会补出原文没有的信息？论文把抽取条目和整合摘要分开检查。三个评判模型报告的条目平均幻觉率为百分之二点三六。
\par
摘要部分，作者移除了具体时间和具体联系等约束。三个模型报告的错误关联率都明显上升，例如GPT的结果从零点六一升到七点四五。这说明来源时间和细节要求会影响摘要质量。条目幻觉和错误关联的统计单位不同，这里主要关注有无约束时的变化。}
\end{frame}

\source{原文 §3--5、Limitations；末项为阅读后提出的实验方向}
\begin{frame}{阅读结论与后续问题}
\small
\begin{columns}[T,onlytextwidth]
\column{.48\textwidth}
\textbf{方法贡献}\par\vspace{2mm}
事实与互动关系共同保存，时间锚支持事件重建，跨事件摘要组织相关历史。\par
\column{.48\textwidth}
\textbf{实验支持}\par\vspace{2mm}
总分76.82\%；整合模块主要改善多跳与时间任务；相对图记忆降低构建开销。\par
\end{columns}
\vspace{4mm}
\textbf{值得继续研究的问题}\par\vspace{2mm}
计划或关系发生变化后，旧条目与摘要怎样同步更新？\par\vspace{1mm}
可用带时间变更的对话，比较仅优先新记忆与显式修订的回答质量、来源可追溯性。\par
\vspace{4mm}
\key{我的理解：记忆写入时的组织方式，会直接影响之后的证据联合。}\par
\vspace{3mm}
{\scriptsize 论文：\url{https://aclanthology.org/2026.acl-short.12/}\par
代码：\url{https://github.com/zjunlp/LightMem/blob/main/StructMem.md}}
\note{建议35秒。\par
最后总结一下。StructMem通过事实与关系抽取保留事件语境，用时间锚恢复相关条目，再用摘要组织跨事件联系。实验支持它对多跳和时间任务的帮助，也展示了相对图记忆的成本优势。
\par
我最关注的后续问题是记忆更新。计划变化以后，旧条目和摘要需要怎样同步修订，才能持续支持可靠回答？这可以作为后续的小规模实验方向。通过这篇论文，我的主要理解是：记忆在写入时怎样组织，会直接影响模型之后怎样联合证据。}
\end{frame}
